Federated learning lets IoT gateways train DDoS detectors without pooling traffic, yet a recent lightweight pipeline fits its feature selection and scaling on pooled, pre-split data, which contradicts data locality. We audit that pipeline, reproduce it on real CIC-IoT2023 flows, and propose FedStat-FS: variance filtering, Pearson and LDA selection, their intersection and global standardisation are computed exactly from two additive uploads per client, and a linear SVM is trained with warm-start FedAvg. On CIC-IoT2023 and N-BaIoT (\nv{seeds} seeds, identically tuned baselines, test data used once) the federated subset equalled the pooled one in \nv{exact_ok} of \nv{exact_total} runs. Under Dirichlet(0.1) label skew, naive local preprocessing fell to F1 \rv{local_prep_fedavg}{ciciot}{dir0.1}{f1m} while FedStat-FS kept \rv{fedstat}{ciciot}{dir0.1}{f1m}; ablations attribute the gap to global standardisation, not to selection. Pooled pre-split preprocessing showed no measurable optimism. The Pearson--LDA intersection, however, brought no accuracy benefit and made linear detectors fail in \nv{nfail:fedstat:nbaiot} of \nv{nrun:fedstat:nbaiot} N-BaIoT runs, whereas Pearson-only never failed. The statistics exchange costs \nv{statbytes:cic}\,KB per client, once. A re-implementation of the base paper's Chimp-optimisation and federated-CNN pipeline was more accurate in every regime (F1 $\ge$ \nv{base:f1min}) but needs pooled raw rows for its feature search and \nv{base:paramratio}$\times$ more parameters.
